\documentclass[12pt]{article}
\usepackage[a4paper,margin=2.5cm]{geometry}
\usepackage[T1]{fontenc}
\usepackage[utf8]{inputenc}
\usepackage{lmodern}
\usepackage{microtype}
\usepackage{needspace}
\usepackage{booktabs}
\usepackage{amsmath}
\usepackage{xcolor}
\usepackage{listings}
\usepackage[hidelinks]{hyperref}
\definecolor{codegray}{RGB}{245,245,245}
\lstset{basicstyle=\ttfamily\small,backgroundcolor=\color{codegray},frame=single,
  breaklines=true,showstringspaces=false,columns=fullflexible,keepspaces=true}
\setlength{\parindent}{0pt}
\setlength{\parskip}{0.55em}

\begin{document}
\pagenumbering{gobble}
\begin{titlepage}
\centering
{\large Ministerul Educa\c{t}iei \c{s}i Cercet\u{a}rii al Republicii Moldova\par}
\vspace{0.45cm}
{\large Universitatea Tehnic\u{a} a Moldovei\par}
\vspace{0.45cm}
{\large Facultatea Calculatoare, Informatic\u{a} \c{s}i Microelectronic\u{a}\par}
\vspace{0.45cm}
{\large Departamentul Ingineria Software \c{s}i Automatic\u{a}\par}
\vspace{3cm}
{\Large\bfseries Report on Laboratory Work No. 12\par}
\vspace{0.55cm}
{\LARGE\bfseries Producer, Consumer and the Other Tools\par}
\vspace{0.35cm}
{\large Operating Systems\par}
\vfill
\begin{flushleft}
\textbf{Done by:} st. gr. FAF-241, Ca\c{t} Chiril\\[0.8cm]
\textbf{Checked by:} asist. univ. Rei\c{t}man Patricia
\end{flushleft}
\vfill
{Chi\c{s}in\u{a}u -- 2026\par}
\end{titlepage}
\pagenumbering{roman}
\tableofcontents
\clearpage
\pagenumbering{arabic}

\section{Aim, environment, and deliverables}
Lab 11 used a semaphore as a lock, a signal, or a resource counter.
This lab combines synchronization tools: a bounded producer/consumer FIFO,
a condition variable paired with a predicate, a reader-preference
readers/writers algorithm, and a two-lock deadlock followed by its repair.
All four required parts are complete.

The programs ran on Ubuntu 24.04.5 LTS on aarch64, using GCC 13.3.0 and
POSIX threads. All eight configurations compiled successfully with
\texttt{-std=c11 -Wall -Wextra -Werror -O0 -pthread}. There were
\textbf{27 executions}: three per correct program, three for the
condition-variable \texttt{if} variation, one each for the two sleeping
broken configurations, and ten original deadlock runs without sleeps.

The required files are \texttt{prodcons.c}, \texttt{condvar.c},
\texttt{readers.c}, and the \textbf{fixed} \texttt{deadlock.c}, under
\texttt{lab12/code/}. The companion \texttt{thread\_utils.h} checks pthread
error codes and semaphore failures, including interrupted waits. Workers
are joined before synchronization objects are destroyed. Stable ID arrays
replace integer-to-pointer casts in the readers/writers example.
The required \texttt{report.md} contains all twelve Observe answers.
This PDF and its editable source follow the previous laboratory reports.
Original output, exit codes, build records, and verification are retained
under \path{lab12/report_files/output/}.

Correct versions are the defaults. Separate compiler definitions create
the experiments without leaving the submitted sources broken:
\begin{lstlisting}
# From the repository root, on Ubuntu:
gcc -std=c11 -Wall -Wextra -Werror -O0 -pthread \
  -o /tmp/prodcons lab12/code/prodcons.c
# Variations (use separate output binary names):
# prodcons.c: -DBAD_LOCK_ORDER=1
# condvar.c:  -DUSE_IF=1
# deadlock.c: -DOPPOSITE_ORDER=1
# deadlock.c: -DOPPOSITE_ORDER=1 -DDELAY_SECONDS=0
\end{lstlisting}
I used \texttt{timeout 5} to bound each deadlock experiment, with a
one-second forced-kill fallback in the automated runner. Exit 124 means
the timeout expired, not that the program completed. Output is line
buffered so messages already printed are retained when a broken process
is stopped. No indefinite background process was left running.

\clearpage
\section{Part 1: producer and consumer}
\subsection{A bounded circular FIFO}
The buffer has four slots and carries eight items. \texttt{empty} starts
at 4, \texttt{full} at 0. The mutex protects the buffer, indices, and
occupancy. Each worker waits for a resource permit \emph{before} acquiring
the mutex, updates the circular buffer, unlocks, and posts the opposite
semaphore. The indices wrap with modulo \texttt{SIZE}.
\lstinputlisting[basicstyle=\ttfamily\footnotesize]{output/prodcons-1.stdout}
All three correct runs reached occupancy four, consumed exactly items
1--8 in FIFO order, and finished empty. A different interleaving would
also be valid. This is the same bounded producer/consumer concept found
in pipes, not a claim that a kernel pipe uses these exact POSIX objects.

\subsection{The wrong order creates a deadlock}
I built the altered producer with \texttt{-DBAD\_LOCK\_ORDER=1}:
\lstinputlisting[basicstyle=\ttfamily\footnotesize]{output/prodcons_bad.txt}
After filling the buffer, the producer holds the mutex while waiting for
an empty permit. The consumer needs that mutex to remove an item and post
such a permit. The circular dependency prevents either thread from making
progress. The final source restores wait-before-lock; a timeout is used
only to terminate the deliberate broken experiment.

\clearpage
\section{Part 2: the condition variable}
\subsection{The predicate is the remembered state}
The ordinary integer \texttt{ready} is the condition, protected by
\texttt{lock}. The writer copies the message, sets \texttt{ready = 1},
and signals while holding that mutex. The reader checks the predicate
under the same mutex and waits in a loop:
\begin{lstlisting}[language=C]
pthread_mutex_lock(&lock);
while (!ready) {
    printf("B: not ready, I sleep\n");
    pthread_cond_wait(&changed, &lock);
}
printf("B: I read \"%s\"\n", message);
pthread_mutex_unlock(&lock);
\end{lstlisting}
This excerpt omits the checked-call wrappers. The wait atomically releases
the mutex and waits, then reacquires it before returning. This allows the
writer to acquire the lock while the reader sleeps and avoids a missed
notification between testing the predicate and starting the wait.

All three correct runs produced:
\lstinputlisting{output/condvar-1.stdout}
If the writer finishes before the reader waits, the reader sees
\texttt{ready == 1} and skips waiting. A condition-variable signal is not
a stored permit; \texttt{ready} preserves the state even if nobody was
waiting when the signal occurred.

\subsection{The if variation works here, but is not the general rule}
The separately compiled \texttt{-DUSE\_IF=1} variant replaced the loop
with a single \texttt{if}. All three runs also printed:
\lstinputlisting{output/condvar_if-1.stdout}
No unwanted wakeup was exposed by these simple one-reader/one-writer runs.
A condition wait may nevertheless return spuriously while the predicate
is false. With more consumers, another thread might consume a resource
before a woken waiter reacquires the mutex. The correct pattern therefore
rechecks the actual predicate in a \texttt{while} loop. This experiment
did not force or claim to observe a spurious wakeup. The submitted default
retains the loop.

\clearpage
\section{Part 3: readers and writers}
\subsection{One permit for the whole reader group}
The first reader acquires \texttt{room}; the last reader releases it.
\texttt{rlock} serializes changes to the reader count and the associated
first/last decisions:
\begin{lstlisting}[language=C]
/* Reader entry, while holding rlock */
if (++readers == 1)
    sem_wait_checked(&room);
/* Reader departure, while holding rlock */
if (--readers == 0)
    sem_post(&room);
\end{lstlisting}
Each writer must acquire the same room permit and keep it throughout its
update and simulated work. Other readers join an active reader group
without acquiring room again. The 200 ms reading intervals remain outside
\texttt{rlock}, so readers can genuinely overlap.

One recorded run was:
\lstinputlisting[basicstyle=\ttfamily\footnotesize]{output/readers-1.stdout}
Each of three runs reached three overlapping readers. The retained event
trace showed no writer overlapping a reader or another writer. Both writes
added ten, leaving value 20 and reader count zero. Which writer runs first
is not a general guarantee.

\subsection{Diagnostic safety and writer starvation}
The handout reads \texttt{readers} in a print outside \texttt{rlock} while
other readers may update it, introducing a diagnostic data race. I moved
that print into the protected region. Additional departure records let
the verifier reconstruct complete active intervals, rather than assuming
a label such as ``alone'' proves exclusivity.

This algorithm gives preference to readers. If new readers keep joining
before the last active reader leaves, the count never reaches zero and a
writer can wait indefinitely: \textbf{writer starvation}. Readers may
still make progress, so this is not the circular-wait deadlock of Part 4.
The finite test finishes both writers; it does not establish fairness
under an unending arrival stream.

\clearpage
\section{Part 4: building and removing a deadlock}
\subsection{Opposite lock order}
In the original version, worker one takes A then B, while worker two takes
B then A. Both sleep after acquiring their first lock. The actual run was:
\lstinputlisting{output/deadlock_original.txt}
One holds A and waits for B; two holds B and waits for A. The wait-for
cycle is \(\text{one} \to \text{two} \to \text{one}\).
The source and ownership messages explain the circular wait; the timeout
alone would not establish the cause of a hang.

\subsection{The submitted repair: A before B everywhere}
Both workers now acquire A then B and release B then A. All three fixed
runs completed normally. One run produced:
\lstinputlisting{output/deadlock_fixed-1.stdout}
\begin{lstlisting}
exit code: 0
\end{lstlisting}
Worker two's message now accurately states that it holds A. The handout's
fixed sample retains an outdated ``I have B'' label. Either worker can
acquire A first in another schedule; the essential property is the common
lock order, not which worker starts first.

\subsection{Removing sleeps does not fix the original}
I compiled out both sleeps while retaining opposite lock order and ran
that binary ten times under a five-second timeout:
\begin{center}
\begin{tabular}{lrrrrrrrrrr}
\toprule
Run & 1 & 2 & 3 & 4 & 5 & 6 & 7 & 8 & 9 & 10 \\
\midrule
Exit code & 124 & 0 & 124 & 0 & 0 & 124 & 124 & 124 & 0 & 0 \\
\bottomrule
\end{tabular}
\end{center}
\textbf{Five completed and five timed out.} A worker sometimes acquires both
locks before its peer establishes the cycle. Other schedules still
produce deadlock. This is harder to diagnose, not a repair: successful
runs hide a timing-dependent defect. Ten outcomes are not a measured
failure probability or a portable failure rate.

\clearpage
\section{Handout questions and answers}
\subsection{Part 1: producer and consumer}
\begin{enumerate}
\Needspace{7\baselineskip}
\item \textbf{Question:} The producer is fast, yet ``in buffer'' never goes above 4. Which call stops it?
\par\textbf{Answer:} \texttt{sem\_wait(\&empty)}, through
\texttt{sem\_wait\_checked}, blocks when no free-slot permit remains.
Four initial permits and balanced consumer posts enforce the capacity;
a mutex alone would not prevent overflow.

\Needspace{7\baselineskip}
\item \textbf{Question:} Who waits at the start: the consumer or the producer? Which call makes it wait?
\par\textbf{Answer:} The consumer waits at \texttt{sem\_wait(\&full)}
when the buffer is still empty, because full starts at zero. The producer
initially has four empty permits. If it produces first, a full permit may
already be available when the consumer reaches wait.

\Needspace{8\baselineskip}
\item \textbf{Question:} Lock the producer's mutex before sem\_wait(\&empty). What happens when the buffer is full, and why?
\par\textbf{Answer:} The altered run filled four slots and timed out with
exit 124. The producer holds the mutex and waits for an empty slot; the
consumer needs the mutex to remove an item and post an empty permit.
Each is waiting for progress that the other cannot make.
\end{enumerate}

\subsection{Part 2: the condition variable}
\begin{enumerate}
\Needspace{7\baselineskip}
\item \textbf{Question:} In Lab 11 the semaphore remembered that the message was ready. Here, which variable remembers it?
\par\textbf{Answer:} \texttt{ready}, an integer protected by the mutex.
The condition variable is a place to wait for possible changes, not a
stored readiness permit.

\Needspace{7\baselineskip}
\item \textbf{Question:} Replace while (!ready) with if (!ready), compile, run. Does it still work here?
\par\textbf{Answer:} Yes, all three altered runs read hello after the
writer's announcement. Those runs did not expose an unwanted wakeup and
do not make a one-time test the correct general pattern.

\Needspace{8\baselineskip}
\item \textbf{Question:} Why write while? Give one reason a thread can wake up and still find ready == 0.
\par\textbf{Answer:} A POSIX condition wait can return spuriously before
the predicate becomes true. More generally another consumer may change
state before the waiter reacquires the mutex. The loop rechecks the
predicate under its lock and waits again when necessary.
\end{enumerate}

\Needspace{12\baselineskip}
\subsection{Part 3: readers and writers}
\begin{enumerate}
\Needspace{8\baselineskip}
\item \textbf{Question:} Three readers are inside together, but a writer is always alone. Which two lines of reader() make that possible?
\par\textbf{Answer:} The first-reader \texttt{sem\_wait(\&room)} and
last-reader \texttt{sem\_post(\&room)} calls are guarded by
\texttt{++readers == 1} and \texttt{--readers == 0}, respectively. They
hold room for the whole group. Writers need that same permit, while
later readers can join the active group.

\Needspace{7\baselineskip}
\item \textbf{Question:} Why does readers need its own lock, rlock? What could go wrong without it?
\par\textbf{Answer:} Updates and first/last-reader decisions must be
serialized. Lost updates or inconsistent first/last decisions could
fail to release room, or release it while a reader remains. The printed
snapshot of the count must also be protected.

\Needspace{7\baselineskip}
\item \textbf{Question:} Readers keep arriving without a pause. What happens to writer 2?
\par\textbf{Answer:} It may starve. If the active reader count never
reaches zero, the group never releases room. Reader preference does not
guarantee bounded waiting for a writer, even though this finite test ends.
\end{enumerate}

\subsection{Part 4: deadlock}
\begin{enumerate}
\Needspace{9\baselineskip}
\item \textbf{Question:} In the stuck run, which worker holds which mutex, and which does each wait for? Draw the circle.
\par\textbf{Answer:} One holds A and waits for B; two holds B and waits
for A. The resource cycle is:
\[
\begin{array}{ccc}
\text{one} & \xrightarrow{\text{waits for}} & B \\
\scriptstyle\text{held by}\ \uparrow & & \downarrow\ \scriptstyle\text{held by} \\
A & \xleftarrow{\text{waits for}} & \text{two}
\end{array}
\]

\Needspace{8\baselineskip}
\item \textbf{Question:} Remove both sleeps from the original and run ten times. Does it still always hang? Why is that worse, not better?
\par\textbf{Answer:} No: five completed and five timed out. The lock-order
bug remains, but depends on interleaving. An intermittent failure can
escape testing and be harder to reproduce than an obvious consistent hang.

\Needspace{7\baselineskip}
\item \textbf{Question:} In one sentence, why can taking the locks in the same order never deadlock?
\par\textbf{Answer:} For these two locks, requiring A before B prevents
circular wait because no worker can hold B while waiting to acquire A.
This removes this lock-order deadlock, not every possible unrelated wait.
\end{enumerate}

\clearpage
\section{Verification and conclusion}
\subsection{Results and evidence}
The retained verifier passed \textbf{253 assertions across 27 executions}.
It reconstructs state from events, rather than requiring the sample's
exact scheduling order:
\begin{itemize}
\item Every correct producer/consumer run produced and consumed items
1--8 exactly once in FIFO order. Every occupancy matched the modeled queue,
stayed within 0--4, and ended at zero.
\item The altered producer reached capacity and timed out with 124. The
ordinary producer/consumer runs all exited normally.
\item Both condition-variable configurations read hello after the writer's
announcement in their three runs. No spurious wakeup was claimed or
injected; the reason for the loop follows from condition-wait semantics.
\item Entry and departure events showed multiple readers concurrently,
no reader/writer or writer/writer overlap, accurate reader counts, and
final shared value 20 after two writes.
\item The original sleeping deadlock printed both first-lock ownership
messages and timed out. All three fixed runs joined both workers and
printed done. Each of ten no-sleep original outcomes was preserved,
including both successes and hangs.
\end{itemize}
Raw stdout and stderr agree with the machine-readable records. The
runner treats unexpected errors or timeouts in correct configurations
as failures. Every deliberate hang was bounded and terminated; no thread
was silently left running. Timing records describe observations, not
correctness thresholds or estimated scheduling probabilities.

\subsection{What the experiments establish}
The synchronization design supplies the safety argument; repeated runs
support it but do not exhaust all schedules. In particular, three working
\texttt{if} runs do not disprove spurious wakeups, and two completed writers
do not establish starvation freedom under continuous reader arrivals.
The protected diagnostic snapshot avoids adding an unrelated data race to
the readers/writers program.

The combined examples illustrate distinct responsibilities. Semaphores
count empty/full slots and hold a room for a reader group; a mutex protects
state updates; a condition variable waits for a predicate protected by its
mutex. Consistent lock ordering removes the circular wait between A and B.
Removing sleeps changes only the chance of exposing that cycle, not the
underlying correctness of the original program.
\end{document}
